\documentclass[11pt]{article}

\usepackage[margin=1in]{geometry}
\usepackage[utf8]{inputenc}
\usepackage[T1]{fontenc}
\usepackage{times}
\usepackage{microtype}
\usepackage{amsmath,amssymb,amsthm}
\usepackage{booktabs}
\usepackage{multirow}
\usepackage{graphicx}
\usepackage{xcolor}
\usepackage[numbers,sort&compress]{natbib}
\usepackage{algorithm}
\usepackage{algpseudocode}
\usepackage[colorlinks=true,linkcolor=blue!60!black,citecolor=blue!60!black,urlcolor=blue!60!black]{hyperref}
\usepackage[capitalize,noabbrev]{cleveref}

% Clean theorem environments (optional)
\newtheorem{proposition}{Proposition}
\newtheorem{definition}{Definition}

% Common math helpers
\newcommand{\R}{\mathbb{R}}
\newcommand{\E}{\mathbb{E}}
\newcommand{\norm}[1]{\left\lVert #1 \right\rVert}

\title{\textbf{Paper Title: Concise Mechanism or Empirical Finding}}

\author{
  Author One$^{1}$ \quad Author Two$^{1}$ \\
  $^{1}$Affiliation \\
  \texttt{contact@example.com}
}

\date{}

\begin{document}
\maketitle

\begin{abstract}
One concise paragraph (110--200 words, 5--7 sentences): establish the dominant
paradigm, state the exact structural bottleneck or anomaly, present the proposed
mechanism or empirical test, report 2--3 concrete quantitative results alongside
their compute/memory budget and strongest baseline, and note scope or code availability.
\end{abstract}

\section{Introduction}
\label{sec:intro}

Keep the introduction to 3--5 tight paragraphs (250--450 words). Open directly
with the technical paradigm and its concrete bottleneck---never with generic
throat-clearing~\citep{vaswani2017attention,he2016deep,hu2022lora}. Place your
teaser or architecture figure on page~1 or top of page~2.

\section{Problem Statement and Background}
\label{sec:background}

Define the formal setting, stream/evaluation contract, and notation before
introducing the method. Contrast with the closest prior classes along explicit
quantitative axes (see \cref{tab:axes}).

\begin{table}[t]
\centering
\caption{\textbf{Structural comparison across mechanism classes.}
$n$: horizon, $d$: representation dimension, $r$: update rank.}
\label{tab:axes}
\begin{tabular}{lccc}
\toprule
Method Class & Update Complexity & Persistent Memory & Exact Rollback \\
\midrule
Baseline A & $O(d^2)$ & $O(d^2)$ & No \\
Baseline B & $O(n \cdot d)$ & $O(n \cdot d)$ & Approximate \\
Proposed & $O(r \cdot d)$ & $O(r \cdot d)$ & Yes \\
\bottomrule
\end{tabular}
\end{table}

\section{Method}
\label{sec:method}

Present the method top-down: overall loop first, then each component in
execution order. Define every symbol and tensor dimension immediately.

\begin{equation}
\label{eq:core}
  \Delta W = B A, \qquad B \in \R^{d_{\mathrm{out}} \times r}, \quad A \in \R^{r \times d_{\mathrm{in}}}.
\end{equation}

\section{Experimental Setup and Results}
\label{sec:experiments}

\subsection{Setup}
State datasets/simulators, held-out seeds or streams, matched budgets, optimizer,
hardware, and wall time concisely. Place full hyperparameter tables in
\cref{app:details}.

\subsection{Main Comparison}
Report quality metrics alongside cost/budget columns with paired confidence
intervals over independent streams or seeds.

\subsection{Ablations and Mechanism Analysis}
Vary one factor at a time from a clearly defined \texttt{base} row, leaving
unchanged cells blank, and answer 2--3 diagnostic questions explaining
\emph{why} the mechanism works or where it fails.

\section{Related Work}
\label{sec:related}

Group prior work by mechanism or assumption rather than paper-by-paper summaries,
ending each paragraph with the exact distinction from this work.

\section{Conclusion}
\label{sec:conclusion}

Two to three short paragraphs: restate what was established, state concrete
limitations and open questions, and link reproducible code and artifacts.

\bibliographystyle{plainnat}
\bibliography{references}

\appendix
\section{Experimental Details and Extended Derivations}
\label{app:details}

Include complete hyperparameter tables, full derivations, and per-stream traces.

\end{document}
